\chapter{Literature Review}
\label{ch:literature}

This chapter reviews the foundational literature that informs the design of the MEDF platform. The review is organized into three areas. First, it examines the foundational work of Zhang et al.\ on the Fairness in Design (FID) framework, which forms the direct theoretical basis for this project. Second, it surveys the broader landscape of AI governance frameworks, focusing on the three selected for implementation in MEDF: the EU ALTAI, NIST AI RMF, and Singapore's MGAF. Finally, it delves into the technical literature on Multi-Criteria Decision-Making (MCDM) methods, which provide the algorithmic foundation for MEDF's scoring and analysis capabilities.

\section{Fairness in Design}
\label{sec:fid}

The primary theoretical inspiration for this project is the work of Zhang, Shu, and Yu on the Fairness in Design (FID) framework~\cite{zhang2023fairness}. Their 2023 paper, \enquote{Fairness in Design: A Framework for Facilitating Ethical Artificial Intelligence Designs}, identifies a critical gap in the AI development lifecycle: the lack of a structured methodology for design teams to brainstorm and surface potential fairness issues during the initial design stage. The authors argue that while many tools exist for post-hoc bias detection in models and data, few resources are available to guide the more foundational, upstream process of designing fairness-aware systems from the outset.

To address this, they propose the FID framework, a game-like methodology centered around a set of prompt cards. These cards are designed to facilitate discussions among team members, encouraging them to consider fairness from the perspectives of various stakeholders. The framework distills the complex field of algorithmic fairness into ten major principles, which serve as brainstorming guides. User studies conducted by the authors demonstrated that FID was effective in helping participants make more nuanced decisions about fairness and lowered the barrier to entry for teams without deep pre-existing knowledge in the field. A related work by Zhang~\cite{zhang2022methodology} further elaborates on this, positioning the FID toolkit as a critical intervention to embed ethical considerations early in the AI pipeline. An earlier conference version~\cite{zhang2021fairness} and a follow-up study on explainability~\cite{zhang2022explainable} provide additional context for this line of research.

MEDF builds directly on the philosophy of the FID framework. It takes the core idea of structured, stakeholder-centric ethical deliberation and seeks to operationalize it within a quantitative, computational platform. While FID is a qualitative, manual tool for brainstorming, MEDF is a quantitative, automated tool for evaluation and analysis. It extends the concept by not only facilitating the identification of issues but also by measuring their severity, quantifying the conflicts between stakeholder views, and proposing mathematically optimized resolutions. In this sense, MEDF can be seen as a downstream complement to the FID methodology: where FID helps shape the initial design, MEDF helps evaluate and govern the resulting implementation.

\section{AI Governance Frameworks}
\label{sec:governance}

The MEDF platform integrates and harmonizes three prominent AI governance frameworks, chosen for their distinct geographical origins and complementary perspectives. A comprehensive survey by Jobin, Ienca, and Vayena~\cite{jobin2019landscape} identified a global convergence around five ethical principles (transparency, justice and fairness, non-maleficence, responsibility, and privacy), while also highlighting substantive divergence in how these principles are interpreted and implemented across different jurisdictions. This finding directly motivates MEDF's multi-framework approach.

\subsection{EU Assessment List for Trustworthy AI (ALTAI)}
\label{sec:altai}

The European Union has taken a strong, rights-based approach to AI regulation, culminating in the proposed AI Act~\cite{ec2021aiact}. The High-Level Expert Group on AI (AI HLEG) published the \enquote{Ethics Guidelines for Trustworthy AI}, which outlines seven requirements for achieving Trustworthy AI: (1) Human Agency and Oversight, (2) Technical Robustness and Safety, (3) Privacy and Data Governance, (4) Transparency, (5) Diversity, Non-discrimination and Fairness, (6) Societal and Environmental Well-being, and (7) Accountability~\cite{aihleg2019ethics}. To make these requirements actionable, the AI HLEG developed the Assessment List for Trustworthy AI (ALTAI)~\cite{aihleg2020altai}. ALTAI is a practical checklist intended to help AI developers and deployers self-assess whether the AI system they are building is trustworthy. It breaks down each of the seven requirements into a series of concrete questions and considerations. Smuha~\cite{smuha2019eu} provides a detailed analysis of the EU's approach to these guidelines. MEDF uses the structure of ALTAI as a primary source for its unified dimension model, mapping ALTAI's requirements directly to its own evaluation criteria.

\subsection{NIST AI Risk Management Framework (AI RMF)}
\label{sec:nist}

Developed by the U.S.\ National Institute of Standards and Technology, the AI Risk Management Framework (AI RMF) provides a voluntary, flexible, and sector-agnostic structure for managing the risks associated with AI systems~\cite{nist2023rmf}. The AI RMF is organized around four core functions: Govern, Map, Measure, and Manage. It is designed to be integrated into an organization's broader risk management processes. Unlike the EU's more prescriptive, rights-based approach, the NIST framework is more focused on providing a process for identifying, assessing, and responding to risks throughout the AI lifecycle. It emphasizes the importance of context, continuous monitoring, and creating a culture of risk management. MEDF incorporates the risk-centric perspective of the NIST AI RMF, particularly in its harm assessment module, which seeks to quantify the potential severity of ethical risks.

\subsection{Singapore Model AI Governance Framework (MGAF)}
\label{sec:mgaf}

Singapore's Model AI Governance Framework (MGAF), now in its second edition, offers a pragmatic, industry-focused perspective on AI governance~\cite{pdpc2020mgaf}. It is designed to be a readily implementable guide for organizations that are developing or deploying AI. The framework is built on two guiding principles: (1) organizations should be accountable for the AI they deploy, and (2) decisions made by AI should be explainable, transparent, and fair. The MGAF provides detailed guidance on internal governance structures, risk management in the AI lifecycle, and measures for customer relationship management. Its practical, business-oriented approach provides a valuable complement to the rights-based focus of the EU and the risk-management process focus of NIST. MEDF's emphasis on configurable stakeholder weights and trade-off analysis aligns well with the MGAF's pragmatic recognition that ethical AI involves balancing different commercial and societal interests.

\section{Multi-Criteria Decision-Making (MCDM) Methods}
\label{sec:mcdm}

The algorithmic core of MEDF is built upon established Multi-Criteria Decision-Making (MCDM) methods. MCDM is a sub-discipline of operations research that deals with making decisions in the presence of multiple, often conflicting, criteria~\cite{behzadian2012topsis}. This is directly analogous to the problem of ethical AI evaluation, where a system must be judged against multiple competing values.

\subsection{TOPSIS}
\label{sec:topsis}

The Technique for Order of Preference by Similarity to Ideal Solution (TOPSIS) is a widely used MCDM method~\cite{hwang1981mcdm}. The fundamental principle of TOPSIS is that the chosen alternative should have the shortest geometric distance from the ideal solution and the longest geometric distance from the anti-ideal solution. The \enquote{ideal solution} is a hypothetical alternative that has the best possible value for each criterion, while the \enquote{anti-ideal solution} has the worst possible value. TOPSIS is computationally efficient and its logic is relatively easy to understand. It normalizes the data for each criterion and then calculates a closeness coefficient for each alternative, allowing them to be ranked~\cite{behzadian2012topsis}. MEDF uses TOPSIS as its primary scoring method due to its robustness and its ability to produce a clear, scalar score that represents the overall ethical performance of an AI system.

\subsection{NSGA-II and Pareto Optimization}
\label{sec:nsga2}

When stakeholders have conflicting priorities, there is often no single solution that is optimal for everyone. In such cases, the goal is to find the set of \enquote{Pareto-optimal} solutions. A solution is Pareto-optimal if it is impossible to improve one stakeholder's outcome without worsening another's. The set of all such solutions is known as the Pareto frontier. NSGA-II (Non-dominated Sorting Genetic Algorithm II) is a popular and effective multi-objective evolutionary algorithm used to find this Pareto frontier~\cite{deb2002nsga2}. It uses concepts of non-dominated sorting and crowding distance to guide its search, allowing it to find a diverse set of high-quality solutions along the frontier. MEDF employs NSGA-II in its Pareto router to generate a set of consensus weight vectors that represent the best possible compromises between conflicting stakeholder priorities.
